\documentclass{beamer}

\usepackage[utf8]{inputenc}
\usepackage[english]{babel}
\usepackage{textpos}
\usepackage{graphicx}
\usepackage{textcomp}
\usepackage{animate}
\usepackage{comment}
%\usepackage{multicol}
%\setlength{\columnsep}{5cm}
\usepackage{amsmath}
\usepackage{amsfonts}
\usepackage{verbatim}
\usepackage{varwidth}
\usepackage{caption}
\usepackage{dcolumn}
\usepackage{amssymb}
\usepackage{algorithm}
\usepackage{algorithmic}
%\usepackage{physics}
\usepackage{hyperref}
\usepackage{scalerel,stackengine}
%\usepackage{fancyvrb}
\usepackage{enumitem}

\usefonttheme[onlymath]{serif}

\newcommand{\code}[1]{\texttt{#1}}
\newcommand{\shp}{$>>>$ }
\newcommand{\tab}{\hspace*{22pt}}

\usetheme{Madrid}

\title{NumPy and Matplotlib}
\author{EEC 3150}
\date{}

\institute{Trevecca Nazarene University}
 


%%%%%%%%% BEGIN DOCUMENT %%%%%%%%%
%%%%%%%%% BEGIN DOCUMENT %%%%%%%%%
%%%%%%%%% BEGIN DOCUMENT %%%%%%%%%

\begin{document}

% get title page 
\frame{\titlepage}
 
%%%%%%%%%%
%%%%%%%%%%
\begin{frame}	
	\frametitle{Outline}
	\tableofcontents	
\end{frame}


%%%%%%%%%%
%%%%%%%%%%
\section{RNG}

\begin{frame}	
	\center{\Huge{RNG}}
\end{frame}

\subsection{Random Walks}

\begin{frame}	
	\frametitle{Random Walks}
	
	\begin{block}{Random Walks:}
		A random walk is a sequence $x_0,x_1,x_2, \cdots, x_{N-1}$ of randomly generated steps through some space:
		\begin{itemize}
			\item It is usually formulated
			\begin{equation}
				x_{i+i} = x_i + r_i
			\end{equation}
			where $i=0, \cdots, N-1$, $x_0$ is a given initial state, and the $r_i$'s are randomly generated from some distribution
		\end{itemize}
	\end{block}
	
\end{frame}

\begin{frame}	
	\frametitle{Random Walks}
	
	\begin{exampleblock}{Create some random walks:}
		\begin{itemize}
			\item 1D +1 or -1
			\item 1D uniform(0,1)
			\item 1D uniform(-1,1)
			\item 1D normal distribution
			\item 2D normal
			\item 2D up-down/left-right (with bias)
		\end{itemize}
	\end{exampleblock}
	
\end{frame}


\subsection{RNG on Different Domains}

\begin{frame}	
	\frametitle{RNG on Different Domains}
	
	\begin{block}{RNG on Different Domains:}
		We can use NumPy's RNG tools to randomly generate points inside an arbitrary domain:
		\begin{itemize}
			\item We can uniformly sample different domains by sampling along the ``dimensions'' of the domain
		\end{itemize}
	\end{block}
	
\end{frame}

\begin{frame}	
	\frametitle{RNG on Different Domains}
	
	\begin{exampleblock}{Uniformly sample some domains:}
		\begin{itemize}
			\item Circle perimeter
			\item Rectangle
			\item Disk area
			\item Sphere
		\end{itemize}
	\end{exampleblock}
	
\end{frame}


%%%%%%%%%%
%%%%%%%%%%
\section{Monte Carlo Simulation}

\begin{frame}	
	\center{\Huge{Monte Carlo Simulation}}
\end{frame}

\begin{frame}	
	\frametitle{Monte Carlo Simulation}
	
	\begin{block}{Monte Carlo Simulation:}
		\textbf{Monte Carlo simulation} is a computational technique used to model and simulate complex systems or processes using random sampling and statistical methods
		
		\vspace{10pt}
		You can use it to:
		\begin{itemize}
			\item Evaluate complex integrals over complex domains
			\item Evaluate statistical measures on complex distributions
			\item Generate complex distributions
		\end{itemize}
		
		\vspace{10pt}
		Technique:
		\begin{itemize}
			\item Generate LARGE volume of random samples on some distribution
			\item Perform various stiatistical analyses of the result
		\end{itemize}
	\end{block}
	
\end{frame}

\begin{frame}	
	\frametitle{Monte Carlo Simulation}
	
	\begin{exampleblock}{Estimate $\pi$:}
		\small
		Steps:
		\begin{enumerate}[label={\arabic*.}]
			\small
			\item Uniformly sample $N_{sq}$ points in the unit square
			\item Inscribe the unit circle in the unit square
			\item Count the number $N_{c\ell}$ points in the unit circle
			\item For large $N_{sq}$, \quad $\frac{N_{c\ell}}{N_{sq}}\approx\frac{A_{c\ell}}{A_{sq}}$
			\item $\frac{A_{c\ell}}{A_{sq}}=\frac{\pi r^2}{(2r)^2}=\frac{\pi}{4} \approx \frac{N_{c\ell}}{N_{sq}}$
			\item $ \pi \approx 4\frac{N_{c\ell}}{N_{sq}}$
		\end{enumerate}
		
		Some notes:
		\begin{itemize}
			\small
			\item More accurate approximations for large $N_{sq}$
			\item MC has a slow rate of convergence, i.e., it takes MANY samples to get accurate results (there are more efficient means of estimating $\pi$)
		\end{itemize}
	\end{exampleblock}
	
\end{frame}

\begin{frame}	
	\frametitle{Monte Carlo Simulation}
	
	\begin{exampleblock}{Estimate $\pi$:}
		\small
		Steps:
		\begin{enumerate}[label={\arabic*.}]
			\small
			\item Uniformly sample $N_{sq}$ points in the unit square
			\item Inscribe the unit circle in the unit square
			\item Count the number $N_{c\ell}$ points in the unit circle
			\item For large $N_{sq}$, \quad $\frac{N_{c\ell}}{N_{sq}}\approx\frac{A_{c\ell}}{A_{sq}}$
			\item $\frac{A_{c\ell}}{A_{sq}}=\frac{\pi r^2}{(2r)^2}=\frac{\pi}{4} \approx \frac{N_{c\ell}}{N_{sq}}$
			\item $ \pi \approx 4\frac{N_{c\ell}}{N_{sq}}$
		\end{enumerate}
		
		Some notes:
		\begin{itemize}
			\small
			\item More accurate approximations for large $N_{sq}$
			\item MC has a slow rate of convergence, i.e., it takes MANY samples to get accurate results (there are more efficient means of estimating $\pi$)
		\end{itemize}
	\end{exampleblock}
	
\end{frame}


\end{document}




